\section{Pilot Decisions}
\label{sec:pilots}

The analysis plan left a few values to pilot data. Each was fixed before
the first campaign from a pilot run on the measuring machine, and is listed
here with the number behind it (Table~\ref{tab:pilots}). Pilot numbers are
evidence for a decision; none is a result, and none is mixed with the
campaigns. The pilots ran on the first build.

\begin{table}[htbp]
\caption{The values fixed from pilots, and the evidence for each. All
numbers in this table are pilot numbers, not results.}
\label{tab:pilots}
\centering\footnotesize
\begin{tabular}{p{0.2\linewidth}p{0.2\linewidth}p{0.5\linewidth}}
\toprule
Decision & Value & Evidence from the pilot \\
\midrule
Image size & 8192$\times$8192 & Smallest size at which the fastest browser cell at $N=8$ is over 100 clock ticks in both browsers (Firefox 213 ticks, Chrome 802) \\
Clock factor & 100 ticks & A factor of 1000 would need images above 20\,000$^2$ in Firefox; 100 caps the quantization of one sample at 1\% \\
Warm-up $W$ & 5 to 71, per workload & 1.5 times the largest warm-up observed in any environment, thread count and run of 100-iteration runs \\
Thread counts & 1, 2, 4, 6, 7, 8 & $N=7$ locates where a knee at the last thread would start \\
Mask & 8 P-cores & With the E-cores added, $N=8$ runs were slower in 5 of 6 cells (3 to 19\%) and several times more variable \\
Display & headless & Headful over headless within 0.3\% in four cells, intervals within 1.2\% \\
Start delay & 3 s & Without it, Firefox's local histogram at $N=8$ was 1.32 times slower (ratio 0.756 [0.721, 0.790]) with bursts in 3 of 10 runs \\
Runs & 30 per cell & 99.5\% half-widths under 1\% for RQ1 and 1.1 to 7.5\% for RQ2 \\
Iterations & 10 per run & More iterations do not improve the agreement of the per-run median (93 to 96\% within 2\% for every count tried) \\
Foreign CPU threshold & 0.25 s & Twice the largest value of 1638 pilot runs (0.125 s) \\
Chrome build & Chrome for Testing & Within 1.3\% of the branded stable Chrome of the same version in four cells \\
Padding & 128 bytes & No difference detected against 64 bytes (Section~\ref{sec:workloads}) \\
\bottomrule
\end{tabular}
\end{table}

\paragraph{Image size and the clock.} A browser clock that is coarse
compared with an iteration makes the measured time mostly quantization
error. The size sweep measured the resolution of \code{performance.now()}
in every browser run, 5~\textmu s in Chrome and 20~\textmu s in Firefox, and
the fastest configuration at $N=8$ in clock ticks: 24, 68, 213 and 322 in
Firefox at sides of 2048, 4096, 8192 and 10\,240. The plan started with a
factor of 1000, which is out of reach at any usable size; the factor became
100, a warning and not an exclusion, and the image 8192$\times$8192, the
smallest size that clears it with a margin of two in Firefox. The sweep
also showed, with 3 runs and no intervals, that the gap between the
environments grows as the image shrinks (filter speedup at $N=8$ of 4.91
natively against 3.45 in Chrome at 2048$^2$, 6.68 against 6.61 at
8192$^2$), which is what the cutoff campaigns later measured properly
(Section~\ref{sec:cutoff}).

\paragraph{Warm-up.} Every workload ran with no warm-up and 100 measured
iterations at $N=1$ and 8. The first iteration is the most expensive
everywhere, natively too (the filter's first iteration takes five times the
steady state at $N=8$: the first write to the output buffer). Most cells
are steady after two or three iterations. The rule: a run's observed
warm-up is the first iteration whose next window of 10 has its median
within a tolerance of the steady state, the tolerance being the larger of
2\% and 1.5 times the largest window deviation inside the steady state, so
that noise is not read as warm-up; a workload's $W$ is 1.5 times the
largest observed value over environments, thread counts and runs, at least
5, and \emph{the same in every environment}. The values: 15 for the filter;
5 for the atomic histogram at every $k$; 71, 21, 15 and 15 for the local
histogram at $k=1$, 4, 16 and 256. The 71 comes from Firefox's local
histogram at $k=1$, which runs about 15\% slower for its first 26 to 47
iterations at $N=8$. A single $W$ of 71 for everything would have cost
about 20 more hours, spent on atomic iterations that show no warm-up.

\paragraph{Thread counts and oversubscription.} A browser runs threads of
its own. At $N=8$ the workers plus those threads may ask for more than the
eight cores of the mask, and the slowdown would look like poor scaling. The
first thread pilot showed a knee in Firefox's local histogram at $N=7$ and
8. The start-delay pilot then removed it: with a 3~s delay no browser shows
a knee. The knee most likely came from Firefox's start-up work competing
with eight workers, not from its threads during the measurement.
Oversubscription stays a possible, much smaller threat, and $N=7$ stays in
the design to show where a knee would start.

\paragraph{The mask, and the first pitfall.} The pilot that compared the
P-core mask with all cores found a problem before it found a result: a
sampler of every thread's allowed CPUs showed Chrome's browser process
outside the mask (Section~\ref{sec:problems}). With that fixed, runs at
$N=1$ agree within 1\% between the masks, and at $N=8$ the wider mask is
slower in 5 of 6 cells, by 3\% (Chrome filter) to 19\% (native local
histogram), with run-to-run variation growing from 0.9\% to 7.2\% for the
native filter: the workers share the cores with the system and land on
slower cores.

\paragraph{Headless.} The rule was fixed before measuring: run headless if
the 95\% interval of headful over headless includes 1 and the ratio is
within 2\%. In the filter and the local histogram at $N=8$, the four ratios
are 0.998 to 1.003 and every interval includes 1.

\paragraph{Start delay.} A freshly launched browser keeps doing start-up
work after the page loads, inside the same mask as the measurement. With no
delay, Firefox's local histogram at $k=16$ and $N=8$ took 5.89~ms against
4.46~ms with a 3~s delay, and its within-run spread (slowest over fastest
iteration) fell from 1.37 to 1.03; the other three cells did not move. A
6~s delay changed the medians by at most 1.3\% against 3~s. Native runs
have no delay, a known asymmetry: it affects $T_1$ and $T_N$ of an
environment alike.

\paragraph{Runs and iterations.} Thirty runs of every cell the primary
contrasts need were subsampled to see how the 99.5\% intervals narrow: with
30 runs the half-width is under 1\% for RQ1 and 1.1\% to 7.5\% for RQ2, the
widest being Chrome at $k=1$. A 5\% target everywhere would need about 60
runs; against RQ2 effects of a factor of two, a few percent is small. Two
thresholds proposed before the pilot (a 5\% half-width, a 1\% agreement of
the per-run median) were dropped after seeing its data, and the achieved
values are reported instead.

\paragraph{Chrome for Testing.} The pinned Chrome is Google's build for
testing, chosen because a fixed version can be downloaded later. Against
the branded stable Chrome of the same version, the filter and the atomic
histogram differ by at most 1.3\% (one interval, the filter at $N=8$,
excludes 1: [1.005, 1.020]), and two versions of the testing build differ
by at most 0.6\%. The comparison ran one build after the other, so a 1\%
effect of the build cannot be told from drift.
